\QHeading{第1题　口算。}
\par\noindent


\begin{oralgrid}
\OralItem{1}{25+75}{100}
\OralItem{2}{36+64}{100}
\OralItem{3}{128+72}{200}
\OralItem{4}{58+42}{100}
\OralItem{5}{89+11}{100}
\OralItem{6}{246+54}{300}
\OralItem{7}{135+65}{200}
\OralItem{8}{92+108}{200}\end{oralgrid}
\par

\QHeading{口算思路补充（对应上表题号）}

\Answer{上面列出得数，下面按小题给出不同计算思路。计数单位法适用于小数，整数题可比较凑整、拆分或商不变等方法。}
\MethodTitle{第8小题：$92+108$}

\MethodBlock{口算思路一：凑整再减去多加的数}{}{\begin{aligned}[t]&\phantom{={}}92+108\\
&=92+110-2\\
&=202-2\\
&=200\end{aligned}}


\par\noindent

\QHeading{第2题　在\(\square\)里填上合适的数。}

\PlainItem{\(45+76=76+\){\color{answerred}45}\((129+58)+42=129+(\)}
\par

\PlainItem{{\color{answerred}58}\(+\){\color{answerred}42}\()\)}
\PlainItem{\(39+\){\color{answerred}82}\(=82+39\)\((36+47)+53=36+(\)}
\PlainItem{{\color{answerred}47}\(+\){\color{answerred}53}\()\)}
\PlainItem{\(128+\){\color{answerred}53}\(=53+128\)\(78+(22+69)=(\)}
\PlainItem{{\color{answerred}78}\(+\){\color{answerred}22}\()+69\)}
\par\noindent

\QHeading{第3题　用简便方法计算下面各题。}

\CalcBlock[方法一：简便计算]{1}{\begin{aligned}[t]&\phantom{={}}\color{black}159+378+241\\&\color{answerred}=(159+241)+378\\&\color{answerred}=400+378\\&\color{answerred}=778\end{aligned}}
\par
\MethodBlock{方法二：按运算顺序直接计算（对照）}{本法只作结果核对；本题要求简便计算，前面的简便解法仍保留。}{\begin{aligned}[t]&\phantom{={}}159+378+241\\
&=537+241\\
&=778\end{aligned}}


\CalcBlock[方法一：简便计算]{2}{\begin{aligned}[t]&\phantom{={}}\color{black}554+348+122\\&\color{answerred}=554+(348+122)\\&\color{answerred}=554+470\\&\color{answerred}=1024\end{aligned}}
\MethodBlock{方法二：按运算顺序直接计算（对照）}{本法只作结果核对；本题要求简便计算，前面的简便解法仍保留。}{\begin{aligned}[t]&\phantom{={}}554+348+122\\
&=902+122\\
&=1024\end{aligned}}

\CalcBlock[方法一：简便计算]{3}{\begin{aligned}[t]&\phantom{={}}\color{black}182+126+318+74\\&\color{answerred}=(182+318)+(126+74)\\&\color{answerred}=500+200\\&\color{answerred}=700\end{aligned}}
\MethodBlock{方法二：交换、结合后先凑整}{保留每个数原来的加减号，先把能凑整的部分放在一起。}{\begin{aligned}[t]&\phantom{={}}182+126+318+74\\
&=(182+318)+126+74\\
&=500+126+74\\
&=626+74\\
&=700\end{aligned}}

\MethodBlock{方法三：按运算顺序直接计算（对照）}{本法只作结果核对；本题要求简便计算，前面的简便解法仍保留。}{\begin{aligned}[t]&\phantom{={}}182+126+318+74\\
&=308+318+74\\
&=626+74\\
&=700\end{aligned}}

\CalcBlock[方法一：简便计算]{4}{\begin{aligned}[t]&\phantom{={}}\color{black}23+(159+77)\\&\color{answerred}=(23+77)+159\\&\color{answerred}=100+159\\&\color{answerred}=259\end{aligned}}
\MethodBlock{方法二：去括号后计算}{加号后的括号去掉不变号，减号后的括号去掉要改变各项的加减号。}{\begin{aligned}[t]&\phantom{={}}23+(159+77)\\
&=23+159+77\\
&=182+77\\
&=259\end{aligned}}

\MethodBlock{方法三：按运算顺序直接计算（对照）}{本法只作结果核对；本题要求简便计算，前面的简便解法仍保留。}{\begin{aligned}[t]&\phantom{={}}23+(159+77)\\
&=23+236\\
&=259\end{aligned}}

\CalcBlock{5}{\begin{aligned}[t]&\phantom{={}}\color{black}287+13+36\\&\color{answerred}=(287+13)+36\\&\color{answerred}=300+36\\&\color{answerred}=336\end{aligned}}
\CalcBlock[方法一：简便计算]{6}{\begin{aligned}[t]&\phantom{={}}\color{black}350+195+105+850\\&\color{answerred}=(350+850)+(195+105)\\&\color{answerred}=1200+300\\&\color{answerred}=1500\end{aligned}}
\MethodBlock{方法二：交换、结合后先凑整}{保留每个数原来的加减号，先把能凑整的部分放在一起。}{\begin{aligned}[t]&\phantom{={}}350+195+105+850\\
&=(350+850)+195+105\\
&=1200+195+105\\
&=1395+105\\
&=1500\end{aligned}}

\MethodBlock{方法三：按运算顺序直接计算（对照）}{本法只作结果核对；本题要求简便计算，前面的简便解法仍保留。}{\begin{aligned}[t]&\phantom{={}}350+195+105+850\\
&=545+105+850\\
&=650+850\\
&=1500\end{aligned}}

\par\noindent

\QHeading{第4题　解决问题。}

\Problem{①学校图书馆第一周借出127本书，第二周借出89本书，第三周借出73本书。三周一共借出多少本书？}
\par


\MethodTitle{方法一}
\Answer{\Formula{\begin{aligned}[t]&\phantom{={}}127+89+73\\&=(127+73)+89\\&=200+89\\&=289\text{（本）}\end{aligned}}\par 答：三周一共借出289本书。}

\MethodTitle{方法二：按先后顺序逐次相加}
\Answer{先合计前两部分，再加第三部分；这是与凑整法对照的常规算法。}
\Answer{\Formula{\begin{aligned}[t]&\phantom{={}}127+89+73\\
&=216+73\\
&=289\text{（本）}\end{aligned}}}
\Answer{答：三周一共借出289本书。}
\Problem{②小明周末去爬山，第一段路走了156米，第二段路走了138米，第三段路走了144米。小明一共走了多少米？}

\MethodTitle{方法一}
\Answer{\Formula{\begin{aligned}[t]&\phantom{={}}156+138+144\\&=(156+144)+138\\&=300+138\\&=438\text{（米）}\end{aligned}}\par 答：小明一共走了438米。}

\MethodTitle{方法二：按先后顺序逐次相加}
\Answer{先合计前两部分，再加第三部分；这是与凑整法对照的常规算法。}
\Answer{\Formula{\begin{aligned}[t]&\phantom{={}}156+138+144\\
&=294+144\\
&=438\text{（米）}\end{aligned}}}
\Answer{答：一共走了438米。}
